\section{Related Work}\label{sec:related}

\subsection{Applied category theory in artificial intelligence}
The last decade has seen a growing use of category theory as a compositional language for machine learning and AI. 
Baez and Stay~\cite{BaezStay2011} introduced a categorical framework for quantum protocols that inspired the compositional semantics of the APEIRON Framework. 
Fong and Spivak~\cite{FongSpivak2019} demonstrate how monoidal categories can structure data and algorithms; our inter‑layer functors are directly inspired by their compositionality principles.
The idea of representing an AI architecture as a network of categories linked by functors also appears in the categorical deep learning programme of Gavranović et al.~\cite{Gavranovic2024}, though the APEIRON category is unique in grounding atomicity as a separated object.

\subsection{Zero‑object characterisations of atomicity}
The equivalence between an object being a zero‑object and a form of foundational atomicity has been explored in mathematical logic and topos theory.
Johnstone~\cite{Johnstone2002} discusses the role of zero‑objects in classifying well‑pointed toposes, which our construction echoes.
In a different direction, Coecke and Kissinger~\cite{CoeckeKissinger2017} use zero‑objects to characterise causality in quantum processes; the similarity with our use of zero‑objects for epistemic irreducibility is striking.

\subsection{Multi‑modal and multi‑axial formal reasoning}
The APEIRON framework verifies atomicity across four independent mathematical frameworks, a methodology reminiscent of multi‑trusted auditing in formal methods.
Rushby~\cite{Rushby2002} advocated combining multiple formal methods to increase assurance, a philosophy that directly motivates our multi‑axial approach.
The concept of “convergence” across independent verification backends has also been studied by Klein et al.~\cite{Klein2009} in the context of OS kernel verification; our use of SymPy, Z3, Lean, and Coq follows a similar multi‑prover paradigm.

\subsection{Autonomous knowledge genesis and foundational ontologies}
The ambition to build a non‑anthropocentric intelligence recalls earlier foundationalist projects.
Hutter's AIXI~\cite{Hutter2005} provides a formal theory of optimal reinforcement learning grounded in Solomonoff induction, though it lacks a layered ontology.
Schmidhuber's Gödel Machine~\cite{Schmidhuber2007} proposes self‑referential code rewriting but does not explicitly construct a categorical model.
The APEIRON Framework’s explicit categorical semantics for atomicity and inter‑layer translation offers a novel bridge between these theoretical AI foundations and practical architectural design.

\subsection{Reproducible categorical verification}
Finally, our work is accompanied by a fully reproducible Python verification script, embracing the open‑science values championed by the \emph{Compositionality} community~\cite{FongSpivak2019}. All tables in this paper are machine‑generated and can be regenerated by the reader.